\documentclass{beamer}

\usepackage{./styles/esi-beamer}


\title{Chapitre 10: Les tableaux}
\subtitle{\tiny [DEV1] Cours de d\'eveloppement}
\author{Andrew SASSOYE}
\date{2019-2020}


\begin{document}
	\begin{frame}
		\titlepage
	\end{frame}

	\begin{frame}
		\begin{center}
			\includegraphics[width=8mm]{./styles/images/cc}
			\includegraphics[width=8mm]{./styles/images/by}
			\includegraphics[width=8mm]{./styles/images/nc}
			\includegraphics[width=8mm]{./styles/images/sa}
		\end{center}

		\begin{block}{Licence}
			\small Cette \oe uvre est mise \`a disposition sous licence Creative Common
			Attribution - Pas d\rq Utilisation Commerciale - Partage dans les M\^emes Conditions 4.0 International.
		\end{block}

		\tiny Pour voir une copie de cette licence, visitez
		\href{https://creativecommons.org/licenses/by-nc-sa/4.0/}{https://creativecommons.org/licenses/by-nc-sa/4.0/}
	\end{frame}

    \part{Th\'eorie}
    \frame{\partpage}

    \section{D\'eclaration et cr\'eation}\label{sec:declaration-et-creation}
        \subsection{Introduction}\label{subsec:introduction}
            \begin{frame}
    \begin{block}{Definition}
        En programmation,
        on parle de tableau pour d\'esigner un ensemble d'\'el\'ements de m\^eme type design\'e par un nom unique,
        chaque \'el\'ement \'etant rep\'er\'e par un indice pr\'ecisant sa position au sein de l'ensemble
    \end{block}
\end{frame}
        \subsection{D\'eclaration de tableaux}\label{subsec:declaration-de-tableaux}
            \begin{frame}{Forme de base}
    La déclaration d'un tableau se fait en ajoutant des crochets à côté du type
    \lstinputlisting[firstline=3, lastline=3, firstnumber=3]{java/Declaration.java}

    \begin{exampleblock}{Types}
        En Java, les éléments d'un tableau peuvent être de n'importe quel type primitif.
    \end{exampleblock}

    \lstinputlisting[firstline=4, lastline=4, firstnumber=4]{java/Declaration.java}

\end{frame}

\begin{frame}{Variantes}
    On peut déclarer plusieurs tableaux sur une seule ligne
    \lstinputlisting[firstline=5, lastline=5, firstnumber=5]{java/Declaration.java}
    Nous pouvons aussi déclarer un tableau sur la meme ligne que de simple variables (de même type).
    Dans ce cas, les crochets se mettent derrière le nom de la variable.
    \lstinputlisting[firstline=6, lastline=6, firstnumber=6]{java/Declaration.java}
\end{frame}

\begin{frame}{Elements de type 'objet'}
    Si l'on a défini une classe de type $Point$, nous pouvons également déclarer un tableau de points.

    \lstinputlisting[firstline=7, lastline=7, firstnumber=7, morekeywords=Point]{java/Declaration.java}
\end{frame}
        \subsection{Cr\'eation de tableaux}\label{subsec:creation-de-tableaux}
            \begin{frame}{Création par l'operateur new}
    En Java, les tableaux sont des objets. Pour créer un nouvel objet on utilise l'opérateur $new$ suivi du nom de la classe\\
    \tiny Plus d'informations sur la programmation orienté objet serons données au cours DEV2 \normalsize

    \lstinputlisting[firstline=9, lastline=9, firstnumber=9, morekeywords=Scanner]{java/OperateurNew.java}
    n étant un nombre naturel. Utiliser un nombre négatif déclanchera une exception!

    \begin{alertblock}{Attention}
        Une fois la taille d'un tableau definie, elle ne pourra pas être modifiée!
    \end{alertblock}
\end{frame}

\begin{frame}{Création par l'operateur new}{Code complet}
    \lstinputlisting[morekeywords=Scanner]{java/OperateurNew.java}
\end{frame}

\begin{frame}{Utilisation d'un initialiseur}
    En Java, nous pouvons immédiatement initialiser les valeurs d'un tableau.
    Dans ce cas, Java sait automatiquement quel est la taille du tableau.
    \lstinputlisting[morekeywords=Scanner]{java/Initialiseur.java}

    \begin{exampleblock}{Expression comme valeur}
        Nous pouvons utiliser des expressions comme valeurs.
        Celles-ci serons évalués avant d'être sauvegardées en mémoire.
    \end{exampleblock}

\end{frame}

\begin{frame}{Utilisation d'un initialiseur}{Variante sans initialiseur}
    Voici à quoi ressemblerais le code précendent sans initialiseur.
    \lstinputlisting[morekeywords=Scanner]{java/SansInitialiseur.java}
\end{frame}

    \section{Utilisation}
        \subsection{Acc\`es individuel aux \'el\'ements d'un tableau}\label{subsec:acces-individuel-aux-elements-dun-tableau}
        \subsection{Affectation de tableaux}\label{subsec:affectation-de-tableaux}
        \subsection{La taille d'un tableau: length}\label{subsec:la-taille-dun-tableau}
        \subsection{Boucle for\ldots each}\label{subsec:boucle-foreach}

    \section{Tableau en argument ou en retour}\label{sec:tableau-en-argument-ou-en-retour}

    \section{Tableaux \`a plusieurs indices}\label{sec:tableaux-a-plusieurs-indices}
        \subsection{Pr\'esentation g\'en\'erale}\label{subsec:presentation-generale}
        \subsection{Initialisation}\label{subsec:initialisation}
        \subsection{Exemple}\label{subsec:exemple}
        \subsection{For\ldots each}\label{subsec:for-each}

    \part{Exercices}
    \frame{\partpage}

\end{document}